% Figure: detached bow shock ahead of a blunt re-entry body versus an
% attached shock on a slender sharp body. Shows the standoff distance,
% the hot shock layer, and why a blunt nose spreads the heat load.
\begin{figure}[htbp]
\centering
\begin{tikzpicture}[scale=1.0]
% === Left panel: blunt body with detached bow shock ===
\begin{scope}[xshift=0cm]
  \node[font=\scriptsize, anchor=south] at (0.2,2.6) {blunt body};
  % freestream arrows
  \foreach \y in {-1.6,-0.8,0,0.8,1.6} {
    \draw[engblue, -{Latex[length=1.4mm]}] (-3.2,\y) -- (-2.5,\y);
  }
  \node[engblue, font=\scriptsize, anchor=east] at (-3.2,2.0) {$M_\infty\!\gg\!1$};
  % blunt nose: arc of large radius centred to the right
  \draw[engblack, very thick] (1.4,1.5) arc (55:-55:1.85);
  \draw[engblack, very thick] (1.4,1.5) -- (2.6,1.5);
  \draw[engblack, very thick] (1.4,-1.5) -- (2.6,-1.5);
  % bow shock: curved, standing ahead of the nose
  \draw[engred, very thick] (-0.1,2.3) .. controls (-1.4,0) .. (-0.1,-2.3);
  \node[engred, font=\scriptsize, anchor=south east] at (-0.7,1.7) {bow shock};
  % shock layer shading hint (hot region between shock and body)
  \node[engamber, font=\scriptsize] at (0.05,0) {hot};
  \node[engamber, font=\scriptsize] at (0.05,-0.45) {layer};
  % standoff distance marker
  \draw[enggray, {Latex[length=1.2mm]}-{Latex[length=1.2mm]}] (-1.05,-1.05) -- (-0.45,-1.05);
  \node[enggray, font=\scriptsize, anchor=north] at (-0.75,-1.1) {$\Delta$};
  % nose radius marker
  \draw[enggray, dashed] (-0.45,0) -- (1.25,0);
  \node[enggray, font=\scriptsize, anchor=south] at (0.5,0.02) {$R_n$};
\end{scope}
% === Right panel: slender sharp body with attached shock ===
\begin{scope}[xshift=7.6cm]
  \node[font=\scriptsize, anchor=south] at (0.2,2.6) {sharp body};
  \foreach \y in {-1.6,-0.8,0,0.8,1.6} {
    \draw[engblue, -{Latex[length=1.4mm]}] (-3.2,\y) -- (-2.5,\y);
  }
  % slender wedge nose
  \draw[engblack, very thick] (-1.0,0) -- (2.6,0.9);
  \draw[engblack, very thick] (-1.0,0) -- (2.6,-0.9);
  % attached oblique shocks springing from the tip
  \draw[engred, very thick] (-1.0,0) -- (1.9,2.0);
  \draw[engred, very thick] (-1.0,0) -- (1.9,-2.0);
  \node[engred, font=\scriptsize, anchor=south west] at (1.0,1.4) {attached shock};
  % small standoff (zero) note at tip
  \node[engamber, font=\scriptsize, anchor=west] at (-0.9,0.55) {$\Delta\!\to\!0$};
  \node[enggray, font=\scriptsize, anchor=west] at (-0.9,-0.6) {$R_n\!\to\!0$};
\end{scope}
\end{tikzpicture}
\caption[Bow shock on a blunt body versus an attached shock on a sharp
body]{Hypersonic flow over a blunt re-entry body (left) and a slender
sharp body (right). The blunt body cannot turn the flow through the
large angle its forward face demands, so the shock detaches and stands
off the nose by a distance $\Delta$ that scales with the nose radius
$R_n$ (for air at high Mach, $\Delta/R_n$ is roughly $0.1$ to $0.2$).
The gas in the shock layer between the bow shock and the body is
processed by a near-normal shock at the stagnation streamline, so it is
hot and slow. The blunt nose pays a large pressure drag but spreads the
stagnation heat load over a large radius, and the stagnation heat flux
falls as $1/\sqrt{R_n}$, so a large $R_n$ lowers the peak flux. The
sharp body keeps an attached shock and low drag, but its tip radius
tends to zero, which drives the local heat flux toward values no
material survives. The design lever for re-entry is bluntness: trade
pressure drag, which decelerates the vehicle high in the atmosphere
where the air is thin, for a survivable stagnation heat flux.}
\label{fig:vol03ch13:blunt-body}
\end{figure}
